\documentclass{article}
\usepackage[margin=2cm]{geometry}
\usepackage{amsmath}
\usepackage{graphicx}
\usepackage{booktabs}
\usepackage[utf8]{inputenc}
\usepackage{ragged2e}
\setlength{\emergencystretch}{3em}
\begin{document}
\sloppy

\noindent\hfill\includegraphics[height=1.8em]{assets/figure_header_logo.png}\par\noindent\hrule\vspace{4pt}

since Io is O(6) and hence vanishes in the e > 0 limit. However, some of the one-point
functions will depend on . It may be the case that only certain choices of  correspond
to supersymmetric schemes, but since the final free energy will be independent of  we will
not worry about this choice. While in principle gives us the free energy, its evaluation on
our numerical solutions is complicated by the integration over u in S6p. As such, we will.
take a slightly roundabout approach to the calculation of the free energy, first calculating
its derivative dF/da and then integrating over the UV parameter a. This will allow us.
to circumvent the integration over u. In order to get dF/da, it will first be necessary to
calculate the one-point functions of the dual field theory operators. This is the topic of the
following subsection.


\subsection*{0.0.1One-point functions}


By the usual AdS/CFT dictionary, the one-point functions of the operators dual to the three.
scalar fields and the metric are given by


\begin{equation}
\begin{array}{l l}{{\langle{\cal O}_{\sigma}\rangle=\operatorname*{lim}{\frac{1}{\epsilon\to0}}{\frac{1}{\epsilon^{3}}}{\frac{1}{\sqrt{\gamma}}}{\frac{\delta S_{r e n}}{\delta\sigma}}{\frac{\delta S_{r e n}}{\delta\phi^{0}}}}}&{{}}\\ {{\langle{\cal O}_{\varphi^{3}}\rangle=\operatorname*{lim}{\frac{1}{\epsilon\to0}}{\frac{1}{\sqrt{\gamma}}}{\frac{1}{\delta S_{r e n}}}}}&{{}}\\ {{\cdot,\quad\ddots}}&{{}}&{{}}&{{}}\end{array}
\end{equation}


We may obtain the explicit values of these vacuum expectation values by varying the on-shell
action . The variation of the counterterm action Sct is straightforward. The variation of S6D
gives rise to one piece which vanishes on the equations of motion, as well as a boundary term
which must be accounted for.We find


\begin{equation}
\begin{array}{l}{{\stackrel{\scriptstyle=\infty}{\scriptstyle|x|\leq1/|}}}\\ {{\qquad{\scriptstyle|x|}}}\\ {{\frac{|y|}{\sqrt{1}}}}\\ {{\scriptstyle|x|}}\end{array}
\end{equation}


Evaluating the limits, we get the following one-point functions.


\begin{equation}
\begin{array}{c}{{\sqrt[[O_{\sigma}\big]=\frac{5}{2}e^{f_{k}}\alpha\beta\Omega}\qquad\langle O_{\sigma^{\prime}}\rangle=\frac{3}{2}e^{-f_{k}}\beta-\frac{15}{8}e^{f_{k}}\alpha^{2}\beta\Omega}}\\ {{\langle O_{\sigma}\rangle=\frac{1}{2}\beta}}\end{array\qquad\qquad(3)\end{array}
\end{equation}

\end{document}
